\documentclass[10pt]{article}

\usepackage{amsmath}
\usepackage{amssymb}
\usepackage{amsthm}
\usepackage{graphicx}
\usepackage{parskip}
\usepackage{xcolor}
\usepackage{pagecolor}
\usepackage[margin=1.2in]{geometry}
\usepackage{enumerate}
\usepackage{enumitem}
\usepackage{tikz}
\newcommand*\circled[1]{%
   \tikz[baseline=(C.base)]\node[draw,circle,inner sep=1.2pt,line width=0.2mm,](C) {#1};}
\newcommand*\Myitem{%
   \stepcounter{enumi}\item[\circled{\theenumi}]}

\usepackage[utf8]{inputenc}
\usepackage[english]{babel}

\usepackage{mathtools}
\DeclarePairedDelimiter\bra{\langle}{\rvert}
\DeclarePairedDelimiter\ket{\lvert}{\rangle}
\DeclarePairedDelimiterX\braket[2]{\langle}{\rangle}{#1 \delimsize\vert #2}

\definecolor{thmcolour}{rgb}{0,0,0}
\definecolor{defcolour}{rgb}{0,0,0}
\definecolor{textcolour}{rgb}{0,0,0}
\definecolor{backgroundcolour}{rgb}{1,1,1}

\pagecolor{backgroundcolour}
\color{textcolour}

\newtheoremstyle{custhm}
{%space above
}{%space below
}{%body font
\color{thmcolour}\em
}{%indent amount
-0em
}{%head font
\bfseries\color{thmcolour}
}{%head punct
}{%after head space
1em
}{%head spec
\thmname{#1}\if\relax\detokenize{#2}\relax:\else\thmnumber{ #2}:\fi\if\relax\detokenize{#3}\relax\else\thmnote{ (#3)}\fi
}

\newtheoremstyle{ex}
{%space above
}{%space below
}{%body font
\color{thmcolour}
}{%indent amount
-0em
}{%head font
\bfseries\color{thmcolour}
}{%head punct
}{%after head space
1em
}{%head spec
\thmname{#1}\if\relax\detokenize{#2}\relax:\else\thmnumber{ #2}:\fi\if\relax\detokenize{#3}\relax\else\thmnote{(#3)}\fi
}

\newtheoremstyle{remark}
{%space above
}{%space below
}{% body font
}{%indent amount
-0em
}{%head font
\bfseries
}{%head punct
}{%after head space
1em
}{%head spec
\if\relax\detokenize{#3}\relax\thmname{#1}:\else\thmname{#3}:\fi
}

\newtheoremstyle{numremark}
{%space above
}{%space below
}{% body font
}{%indent amount
-0em
}{%head font
\bfseries
}{%head punct
}{%after head space
1em
}{%head spec
\thmname{#1}\thmnumber{ #2}:
}

\newtheoremstyle{cusdef}
{%space above
}{%space below
}{%body font
\color{defcolour}
}{%indent amount
-0em
}{%head font
\bfseries\color{defcolour}
}{%head punct
}{%after head space
1em
}{%head spec
%if numbered, include number
%if named, include name
\thmname{#1}\if\relax\detokenize{#2}\relax:\else\thmnumber{ #2}:\fi\if\relax\detokenize{#3}\relax\else\thmnote{ (#3)}\fi
}

\theoremstyle{custhm}
\newtheorem{theorem}{Theorem}[section]
\theoremstyle{cusdef}
\newtheorem{defin}[theorem]{Definition}
\theoremstyle{custhm}
\newtheorem{lemma}[theorem]{Lemma}
\theoremstyle{custhm}
\newtheorem{cor}[theorem]{Corollary}

\theoremstyle{custhm}
\newtheorem{prop}[theorem]{Proposition}

\theoremstyle{ex}
\newtheorem{ex}[theorem]{Example}

\theoremstyle{custhm}
\newtheorem*{theorem*}{Theorem}

\theoremstyle{cusdef}
\newtheorem*{defin*}{Definition}

\theoremstyle{remark}
\newtheorem*{remark*}{Remark}

\theoremstyle{remark}
\newtheorem{remark}[theorem]{Remark}

\theoremstyle{numremark}
\newtheorem{numremark}[theorem]{Remark}

\setcounter{section}{-1}

%\marginpar{to describe which lecture it is}

\newcommand{\N}{\mathbb{N}}
\newcommand{\Z}{\mathbb{Z}}
\newcommand{\Q}{\mathbb{Q}}
\newcommand{\R}{\mathbb{R}}
\newcommand{\C}{\mathbb{C}}
\newcommand{\e}{\mathrm{e}}
\newcommand{\ra}{\rightarrow}
\newcommand{\lef}{\left(}
\newcommand{\res}{\right)}
\newcommand{\ie}{\textit{i.e.}}
\newcommand{\eps}{\varepsilon}
\newcommand{\E}{\mathbb{E}}
\newcommand{\suminf}{\sum_{n=0}^{\infty}}
\newcommand{\suminfa}[1]{\sum_{#1=0}^{\infty}}
\renewcommand{\P}{\mathbb{P}}
\newcommand{\undf}[1]{\textit{\textbf{#1}}}
\renewcommand{\L}{\mathcal{L}}
\renewcommand{\it}[1]{\textit{#1}}
\newcommand{\M}{\mathcal{M}}
\renewcommand{\phi}{\varphi}
\newcommand{\proves}{\vdash}
\newcommand{\lra}{\leftrightarrow}

\renewcommand{\bar}{\overline}
\renewcommand{\O}{\mathcal{O}}


\newcommand{\ac}[1]{\mathcal{#1}}
\newcommand{\A}{\mathcal{A}}


\renewcommand{\subset}{\subseteq}

\renewcommand{\th}{\textrm{th}}

\usepackage{graphicx}

\newcommand{\hj}{\textrm{HJ}}

\title{Ramsey Theory}
\author{Lectures by Imre Leader}
\date{}

\begin{document}
\maketitle
\section{Introduction}
Ramsey Theory is all about the following question:

\begin{center}
    \underline{\textbf{Can we find some order in (enough) disorder?}}    
\end{center}

In a sense, the entire course is about answering this question in different settings.

Chapter 1: Monochromatic systems. Abstract and concrete.

Chapter 2: Partition Regular Equations. More concrete; looking at the naturals with addition and multiplication, and asking about order/disorder there.

Chapter 3: Infinite Ramsey Theory. Abstract; taken to a limit (countable only though).

Prerequisites: None (other than some basic concepts of topology \it{e.g.} open/closed/compact sets).

Literature: In theory, this course is self-contained. But if you would like a different viewpoint/some reinforcement, consider:
\begin{itemize}
    \item Bollobas, ``Combinatorics'', C.U.P. 1986 (For Ch. 3)
    \item Graham, Rothschild, Spencer, ``Ramsey Theory'', Wiley 1990 (For Ch. 1\&2)
\end{itemize}

Example Sheets: this a 16 lecture course, so there are 3 sheets and 3 classes.

\clearpage
\section{Chapter 1}
\subsection{Monochromatic Systems}

This is the introductory chapter from which everything else will flow.

In this course, we take $\N = \{1,2,3,\dots\}$, and write $[n] = \{1,2,\dots,n\}$. For a set $X$ and $r\in \N$, we write $X^{(r)} = \{A\subset X : |A| = r\}$; this is the collection of all $r$-sets in $X$.

\underline{Queston}: Suppose that we are given a $2$-colouring of $\N^{(2)}$, \it{i.e.} a $c:\N^{(2)}\ra\{1,2\}$. Can we find an infinite subset $M\subset \N$ such that $M$ is monochrome, \it{i.e.} $c$ is constant on $M^{(2)}$.

Let's try some examples to get a feel for it.

\begin{remark*}[Examples]\ 
    \begin{enumerate}
        \item Colour $ij$ red if $i + j$ is even, and blue if $i+j$ is odd.
        
        Here we can do it, rather easily - take $M = \{2,4,6,8,\dots\}$, or any subset thereof, or all the odd numbers, \it{etc...}.

        \item Colour $ij$ red if $\max\{n:2^n|i+j\}$ even, and blue otherwise.
        
        This is also a yes, \it{e.g.} $M = \{4^0,4^1,4^2,\dots\}$.

        \item Colour $ij$ red if $i+j$ has an even number of (distinct) prime factors, and blue otherwise.
        
        This is also a yes, but it is not obvious how...
    \end{enumerate}
\end{remark*}

\begin{theorem}[Ramsey's Theorem]
    Let $c$ be a $2$-colouring of $\N^{(2)}$. Then there exists an infinite monochromatic $M\subset \N$.
\end{theorem}
\begin{proof}
    Pick any $a_1 \in \N$. There are infinitely many edges out of $a_1$, so infinitely many have the same colour; say all edges from $a_1$ to infinite set $B_1$ have colour $c_1$.

    Now within $B_2$, we take a point $a_2$ and find an finite $B_2\subset B_1\backslash\{a_2\}$ such that all edges from $a_2$ to $B_2$ are the same colour - these may be red or blue.

    Repeat this process again, within $B_3$, and repeat ad infinitum.

    This gives us an infinite sequence $(a_i)_{i=1}^{\infty}$ of points, and infinite sequence of colours $c_i$ such that edge $a_i$ to $a_j$ with $i<j$ has colour $c_i$.

    Take a constant subsequence of $c_i$, say $c_i:i\in I$. Then $M = \{a_{i}:i\in I\}$ is monochromatic.
\end{proof}
\begin{remark*}[Remarks]\ 
    \begin{enumerate}
        \item This is sometimes called a `2-pass' proof, because we had to do the whole induction, and then go over it again to finish it off.
        \item In example 3, no such example is known!
        \item The same proof shows that whenever $\N^{(2)}$ is $k$-coloured (\it{i.e.} $c:\N^{(2)}\ra [k]$), there still exists an infinite monochromatic set.
        
        Alternatively, we could deduce this from Theorem 1 plus induction. We could view the colours as `1' and `2 or 3 or $\dots$ or $k$' and apply Theorem 1; if we get an infinite $1$-set, then done, and if we get an infinite `2 or $\dots$ or $k$'-set then done by induction on $k$.

        \item An infinite monochromatic set is \undf{more} than having arbitrarily large finite monochromatic sets, \it{e.g.} take disjoint blue $K_2$, $K_3$, $K_4$, and so on, then connect everything remaining with red edges. While this (of course) does not contradict Ramsey's Theorem, we clearly have arbitrarily large blue sets (the $K_n$s), but there is no \it{infinite} blue set.
    \end{enumerate}
\end{remark*}

\begin{remark*}[Example]
    Any sequence $x_1,x_2,\dots$ in $\R$ (or any totally ordered set) has a monotone subsequence.

    Indeed, $2$-colour $\N^{(2)}$ by giving $ij$ ($i<j$) colour \it{up} if $x_i < x_j$, and \it{down} if $x_i > x_j$. Then apply Theorem 1.
\end{remark*}

\reversemarginpar{Lecture 2}

What if we coloured $\N^{(r)}$, say for $r = 3,4,\dots$. Given a $2$-colouring $c:\N^{(r)}\ra \{1,2\}$, must there be an infinite monochromatic set?

For instance, $r = 3$: colour $ijk$ ($i < j < k$) red if $i | j+k$, and blue otherwise. Here we can do this easily - just take $M = \{2^0,2^1,2^2,\dots\}$.

\begin{theorem}[Ramsey for $r$-sets]
    Let $r \in \N$. Then whenever $\N^{(r)}$ is $2$-coloured, there exists an infinite monochromatic set.
\end{theorem}
\begin{proof}
    In the previous proof, when we picked $a_1$ and looked at the lines out of it to other points, we in fact used the $r = 1$ case on the colouring induced on the singletons in the neighbourhood of $a_1$. Armed with these ideas, the proof here will act in exactly the same way.

    $r = 1$ is trivial (just infinite pigeonhole), or if you prefer $r = 2$ is Theorem 1. We apply induction on $r$. Thus suppose the result holds for $r - 1$. Given $c:\N^{(r)}\ra \{1,2\}$, we look at the induced colouring.

    Pick $a_1 \in \N$, and look at $(\N\backslash\{a_1\})^{(r-1)}$. This has an induced colouring $d$ given by $d(F) = c(F\cup\{a_1\})$. Now by induction there is an infinite $B_1\subset \N\backslash\{a_1\}$ such that all the $r-1$-sets have the same colour according to $d$, \it{i.e.} $c(F\cup\{a_1\}) = c_1$ for all $F\subset B^{(r-1)}$, for some colour $c_1 \in [r]$.

    Repeating, we have $a_1 \in B_1$, and infinite $B_2\subset B_1\backslash \{a_2\}$ such that all $F\cup \{a_2\}$, $F\in B_2^{(r-1)}$, have the same colour $c_2$.

    We keep going to infinity, giving us an infinite sequence of distinct points $a_1,a_2,\dots$ and colours $c_1,c_2,\dots$ such that $c(a_{i_1},a_{i_2},\dots,a_{i_r})=c_{i_1}$ ($i_1 < i_2 < \dots < i_r$). There is then an infinite index set $I$ such that $c_i:i\in I$ is constant, and so $M = \{a_i:i\in I\}$ is monochrome.
\end{proof}

\begin{remark*}[Example]
    We saw from Theorem 1 that given points $(1,x_1),(2,x_2),\dots$ in $\R^2$, there exists a subsequence such that the induced (piecewise-linear) function is monotone.

    Functions can have other properties that are stricter; for instance, we could ask for the function to be convex/concave - in fact we can ensure that this is the case.

    On the surface, this seems really hard - but not for us, with $r$-set Ramsey.

    Indeed, just $2$-colour $\N^{(3)}$ by giving $ijk$ colour \it{convex} if they form a convex shape, and \it{concave} (otherwise) if they from a concave shape (any three points must fall into one of the two categories).
    
    We get an infinite monochrome subsequence; the induced function is either convex or concave for any of the three points; and so the overall function is convex/concave.
\end{remark*}

How long does the proof take in the $r = 3$ case? Well, to find each $a_i$ we need to do an infinite two-pass proof (of the $r = 2$ case). So this happens an infinite number of times, and then theres another pass at the end. Essentially, it takes a very long time.

Surprisingly, the infinite version of Ramsey \it{implies} the finite version.

\begin{theorem}[Finite Ramsey]
    For all $m, r$ there exists an $n$ such that whenever $[n]^{(r)}$ is $2$-coloured, there exists a monochromatic $m$-set.
\end{theorem}
\begin{proof}
    Suppose not. We will show that there ia a $2$-colouring of $\N^{(r)}$ without a monochromatic $m$-set, (massively) contradicting Theorem 2.

    For each $n\ge r$, have a $2$-colouring $c_n$ of $[n]^{(r)}$ with no monochromatic $m$-set. [We want to take their union to get a bad colouring of the whole of $\N^{(r)}$, but the problem is that the colours aren't necessarily \it{nested}, \it{i.e.} any two agree where both defined.]

    There are only finitely many ways to $2$-colour $[r]^{(r)}$ (two, in fact), so infinitely many of the $c_n$ agree on $[r]^{(r)}$: say $c_n\restriction [r]^{(r)} = d_r$ for all $n\in A_1$, some $d_n:[r]^{(r)}:\{1,2\}$.

    There are only finitely many ways to $2$-colour $[r+1]^{(r)}$, so infintely many of the $c_n:n\in A_1$ agree on $[r+1]^{r+1}$: say $c_n\restriction [r+1]^{(r)} = d_{r+1}$ for all $n \in A_2$, for some $d_{r+1}:[r+1]^{(r)}\ra \{1,2\}$.

    Continue inductively. We obtain colourings $d_n:[n]^{(r)}\ra\{1,2\}$ for each $n\ge r$ such that
    \begin{enumerate}[label=\arabic*)]
        \item $d_n$ has no monochromatic $m$-set, since $d_n = c_n'\restriction [n]^{(r)}$ for some $n' \ge n$
        \item $d_{n'}\restriction [n]^{(r)} = d_n$ for all $n' \ge n$ by construction of the $d_n$s
    \end{enumerate}

    Now put $c(F) = d_n(F)$, any $n\ge \max F$ (for each $r$-set $F$). We can say \it{any}, because all the colourings agree.

    This is well-defined (by 2)), and has no monochromatic $m$-set (by 1)). Massive contradiction.
\end{proof}

\begin{remark*}[Remarks]
    This is called a `compactness' argument, similar to the proof of Bolzano-Weierstrass in IA Analysis. What we are essentially doing is showing that if the space $\{1,2\}^{\N}$ of $2$-colourings, with the product topology (\it{i.e.} induced from the metric $d(f,g) = 1/(\min \{n:f(n)\ne g(n)\})$) is (sequentially) compact.

    \underline{Note}: This proof also gives \it{no} information on how large $n = n(m,r)$ can be (such proofs do exists though, \it{c.f.} IID Graph Theory).
\end{remark*}

What if we colour $\N^{(2)}$ with infinitely many colours, \it{i.e.} have $c:\N^{(2)}\ra X$, some set $X$?

Of course, we do not get an infinite monochromatic set, since \it{e.g.} we can just colour each edge with a unique colour. But we can ask a slightly different question...

Do we always get an infinte ste $m$ such that $c\restriction m^{(2)}$ is either constant or injective?

Sadly, the answer is still no. We can achieve this by colouring $ij$ ($i<j$) with colour $i$. This clearly is neither injective nor constant on any infinite subset.

Another option  here is that we can colour $ij$ with colour $j$ ($i<j$) - then each colour class is finite, instead of infinite. As it happens, at least one of these four situations must arise:

\begin{theorem}
    Let $c$ be a colouring of $\N^{(2)}$ with an arbitrary set of colours. Then there exists an infinite $M\subset \N$ such that one of the following holds:
    \begin{enumerate}[label=\roman*)]
        \item $c$ is constant on $M^{(2)}$
        \item $c$ is injective on $M^{(2)}$
        \item $\forall ij,kl \in M^{(2)}$, $c(ij) = c(kl)\iff i = k$
        \item $\forall ij,kl \in M^{(2)}$, $c(ij) = c(kl)\iff j = l$
    \end{enumerate}
\end{theorem}

[It is worth remarking that this trivially implies Theorem 1: if there are only finitely many colours, then cases ii), iii), iv) cannot happen.]

\begin{proof}
    There is a superbly nice idea here. The act of comparing pairs of edges is the same as asking a question about the 4-set of the vertices, and we can use Ramsey's Theorem for $r = 4$ to help us here.

    Define a $2$-colouring of $\N^{(4)}$ by giving $ijkl$ colour \it{same} if $c(ij) = c(kl)$, and \it{dif} if $c(ij)\ne c(kl)$. Note that, as ever, the notation implies $i<j<k<l$.

    By Ramsey for $4$-sets, ther exists an infinite monochromatic $M_1$ for this colouring. If the colour is \it{same}, then we have case i). Indeed, given $ij,kl\in M_1^{(2)}$, choose $m < n$ in $M_1$ with $m > i,j,k,l$. Then $c(ij) = c(mn)$, and $c(kl) = c(mn)$ (this deals with the case $j= k$, and other anomalies). So we may otherwise assume that $M_1$ has colour diff.

    Now, $2$-colour $M_1^{(4)}$ by giving $ijkl$ colour \it{same} if $c(jk) = c(il)$, and \it{diff} if not.

    Again by $Ramsey$-4, we have infinite $M_2\subset \M_1$ monochromatic for this colouring. Note that we cannot have $M_2$ be the colour \it{same}, as otherwise pick $i < j < k < l < m < n \in M_2$. We would then have $c(jk) = c(in) = c(lm)$, contradicting $j,k,l,m \in M_1$. Thus $M_2$ is colour \it{diff}.

    The final type of edge pairs in a $4$-set could be interlocking:

    $2$-colour $M_2^{(4)}$ by giving $ijkl$ colour \it{same} if $c(ik) = c(jl)$, and \it{diff} if not. Again we obtain an infinite monochromatic $M_3 \subset M_2$. Once again, we cannot have $M_3$ colour \it{same}, else we choose $i < j < k < l < m < n \in M_3$, and we have $c(im) = c(kn) = c(jl)$, contradicing the fact that $i,j,l,m \in M_2$. Thus $M_3$ is colour \it{diff}.

    So now we know that any two edges in $M_3$ have different colours if they are not adjacent; we now deal with the adjacent case.

    $2$-colour $M_3^{(3)}$ by giving $ijk$ colour \it{same} if $c(ij) = c(jk)$, and \it{diff} if not. Have an infinite monochromatic $M_4 \subset M_3$. We cannot have $M_4$ colour \it{same}, else pick $i < j < k < l \in M_4$, and we have $c(ij) = c(jk) = c(kl)$, contradicting the above. So $M_4$ is colour \it{diff}.

    There are other ways that edge pairs can be adjacent, and these account for the various cases in the Theorem.

    $2$-colour $M_4^{(3)}$ by giving $ijk$ colour \it{left-same} if $c(ij) = c(ik)$, and \it{left-diff} if not. We obtain infinite monochromatic $M_5\subset M_4$.

    Finally, $2$-colour $M_5^{(3)}$ by giving $ijk$ colour \it{right-same} if $c(jk) = c(ik)$, and \it{right-diff} otherwise.

    We obtain infinite monochromatic $M_6\subset M_5$.
\begin{center}
    \includegraphics[scale=0.8]{together.jpg}
\end{center}

    \underline{If $M_6$ is \it{left-diff} and \it{right-diff}}: All pairs of edges have a different colour; this is case 2.

    \underline{If $M_6$ is \it{left-same} and \it{right-diff}}: If edges meet at the left they are the same; at the right they are different; this is case 3.

    \underline{If $M_6$ is \it{left-diff} and \it{right-same}}: Similarly, this is case 4.

    \underline{If $M_6$ is \it{left-same} and \it{right-same}}: Pick $i < j < k \in M_6$. Then $c(ij) = c(ik) = c(jk)$, contradiction.
\end{proof}

\begin{remark*}[Remarks]\ 
    \begin{enumerate}
        \item We could use just one colouring, by colouring a $4$-set $ijkl$ with the partition of $[4]^{(2)}$ induced by $c$ on $\{i,j,k,l\}$. The number of colours is then the number such partitions, which is ${4\choose 2} = 6$. We didn't do this because it seems a bit magical and out of the blue, and slightly obscures what's going on. There are also lots of symbols, making it a bit unpleasant.
        
        \item Similarly, if $c:\N^{(r)}\ra X$ is an arbitrary colouring, we get an infinite $M\subset \N$ and a set $I\subset [r]$ such that $\forall x_1,\dots,x_r,y_1,\dots,y_r\in M^{(r)}$ we have $c(x_1,\dots,x_r) = c(y_1,\dots,y_r) \iff x_i=y_i\ \forall i\in I$, where $I$ is an index set. These $2^r$ coloursing are called the \undf{canonical} colourings of the $r$-sets.
        
        For instance, $ = 2$:
        \begin{itemize}
            \item Case i) corresponds to $I = \emptyset$
            \item Case ii) corersponds to $I = \{1,2\}$
            \item Case iii) corresponds to $I = \{1\}$
            \item Case iv) corresponds to $I = \{2\}$
        \end{itemize}
    \end{enumerate}    
\end{remark*}

\subsection{Van der Waerden's Theorem}

If we $2$-colour $\N$, can we find $3$ consecutive poitns of the same colour?

\underline{Answer}: of course not; just colour $\N$ alternately.

What about $3$ equally spaced points, \it{i.e.} $(a,a+d,a+2d)$? This is not obviously false. If it were true though, that would be nice - we would have found some order amongst the disorder. Could we find more, perhaps $4$ points, or even $m$?

\underline{Aim}: For every $m$, whenever $\N$ is $2$-coloured there exists a monochromatic artihmetic progression of length $m$.

Just to be clear, by length we mean the number of terms \it{i.e.} a sequence $\{a, a + 2,\dots, a+(m-1)d\}$ has length $m$.

This is Van der Waerden's Theorem, and it is very hard to solve.

By our usual compactness argument, this is the same as:

\underline{Aim'}: $(\forall m)(\exists N)$ such that whenever $[n]$ is $2$-coloured, there exists a monochromatic AP of length $m$.

Indeed, if this is false then there is an $m$ such that for every $n \ge m$ there is a colouring $c_n$ of $[n]$ with no monochromatic AP of length $m$. We want to combine these into one big colouring of $\N$, but we can't yet.

But infinitely many $c_n$ agree on $[m]$, and, of those, infinitely many agree on $[m+1]$, and so on.... Put together those (nested) restrictions to obtain a $2$-colouring of $\N$ with no monochromatic AP of length $m$. Contradiction.

In proving this (Aim'), one key idea is to generalise: we in fact show that $\forall m,k\ \exists n$ such that whenever $[n]$ is $k$-coloured there exists a monochromatic AP of length $m$.

\underline{Note}: proving a \it{stronger} result might be \it{easier}, \it{e.g.} in a proof by induction.

Another key idea: given APs $A_1,\dots,A_r$, each of length $m$ - so $A_i = \{a_i,a_i+d_i,\dots,a_i+(m-1)d_i\}$ - we say they are \undf{focused} at $f$ if $a_i + md_i = f$ for each $i$. \it{E.g.} $\{1,4\}$ and $\{5,6\}$ are focused at $7$.

If in addition each $A_i$ is monochromatic (for a given colouring), with no two the same colour, we say they are \undf{colour-focused}. Why do we care?

In a $k$-colouring, if we have APs $A_1,\dots,A_k$, each of length $m-1$, that are colour-focused, then we actually have a monochromatic AP of length $m$, by asking ``what colour is the focus?''

Write $W(m,k)$ for the least $n$ (if it exists) such that whenever $n$ is $k$-coloured, there exists a monochromatic AP of length $m$.

\begin{prop}
    $\forall k, \exists n$ such that whenever $[n]$ is $k$-coloured there exists a monochromatic AP of length $3$.
\end{prop}
\underline{Note}: This will be contained in Theorem 6; it is included here for clarity.


\begin{proof}
    \underline{Claim}: $\forall r \le k, \exists n$ such that whenever $[n]$ is $k$-coloured, we have either:
    \begin{itemize}
        \item a monochromatic AP of length 3, or
        \item $r$ colour-focused APs of length 2
    \end{itemize}
    Given this: put $r = k$ and look at the focus. Whatever colour it is, we get a monochromatic AP of length 3.
    
    \underline{Proof of Claim}: Induction on $r$. $r = 1$ is trivial (can take $n = k+1$).

    Given $n$ suitable for $r - 1$, we'll show that $(k^{2n}+1)\cdot 2n$ is suitable for $r$. So given a $k$-colouring of $[(k^{2n}+1)2n]$, with no mono AP of length 3:

    Break up $[(k^{2n}+1)2k]$ into blocks of length $2n$, say $B_1,B_2,\dots,B_{k^{2n}+1}$, where $B_i = [(i-1)2n+1,i2n]$ (square brackets denote interval with endpoints).

    Inside any interval of length $2n$, we have $r-1$ colour-focused APs of length 2 (by choice of $n$), together with their focus (as length = $2n$).

    Now, the number of ways to $k$-colour a block is $k^{2n}$, and since we have $k^{2n}+1$ it must be the case that some two are identically coloured, say $B_s$ and $B_{s+t}$. Inside $B_s$, we have $r-1$ colour-focused APs of length $2$, say $\{a_1,a_1+d_1\},\dots,\{a_{r-1},a_{r-1}+d_{r-1}\}$ focused at $f$.

    But now the APs $\{a_1,a_1+d_1+2nt\},\dots,\{a_{r-1},a_{r-1}+d_{r-1}+2nt\}$ are colour-focused at $f + 4nt$, and $\{f,f+2nt\}$ is also focused there, giving $r$ colour-focused APs of length $2$. So we have finished the induction, so proved the claim, so finished the proof.
\end{proof}
\begin{remark*}[Remarks]\ 
    \begin{enumerate}
        \item The idea of looking at the number of ways to colour a block is called a \undf{product argument}.
        
        \item The proof shows that $$W(3,k) \le \left. k^{k^{.^{.^{k^{4k}}}}}\right\rbrace k$$
        
        So \it{e.g.} $W(3,3)\le 3^{3^{12}}$. This is a `tower-type' bound.
    \end{enumerate}
\end{remark*}

We are now better-equipped to tackle the full theorem:

\begin{theorem}[Van der Waerden's Theorem]
    $\forall m,k,\exists n$ such that whenever $[n]$ is $k$-coloured there exists a monochromatic AP of length $m$.
\end{theorem}
\begin{proof}
    We induct on $m$. $m=1$ is trival (or $m = 2$ is pigeonhole, or $m = 3$ is Prop 1.5).

    So we may assume that that $W(m-1,k)$ exists for every $k$.

    \underline{Claim}: for every $r\le k$, there exists $n$ such that whenever $[n]$ is $k$-coloured, we have either:
    \begin{itemize}
        \item a monochromatic AP of length $m$, or
        \item $r$ colour-focused APs of length $m-1$
    \end{itemize}
    [Given this, put $ r = k$ and look at the focus to get a mono AP of length $m$.]

    \underline{Proof of Claim}: Induction on $r$: $r = 1$ (take $n = W(m-1,k)$).
    Given $n$ suitable for $r-1$, we'll show that $W(m-1,k^{2n})\cdot 2n$ is suitable for $r$.

    So, we are given a $k$-colouring of $[W(k-1,k^{2n})2n]$, with no mono AP of length $m$:

    Break up $[W(m-1,k^{2n})]$ into blocks of length $2n$, say $B_1,B_2,\dots,B_{W(m-1,k^{2n})}$, where $B_i = [(i-1)2n+1,i2n]$.

    As in the proof before, we need three equally spaced identical blocks - but this is much harder to get.

    The number of ways to $k$-colour a block is $k^{2n}$, so since we have $W(m-1,k^{2n})$ \it{blocks}, we must have (by definition of $W(m-1,k^{2n})$) some  $m-1$ equally spaced blockes that are coloured identically. Say $B_s, B_{s+t},\dots,B_{s+(m-2)t}$.

    Now, inside $B_s$ we have $r-1$ colour-focused APs of length $m-1$ (by definition of $n$), together with their focus (as length = $2n$): say $A_1,\dots,A_{r-1}$ focused at $f$, where $A_i$ has first term $a_i$ and common difference $d_i$. Then the APs $A_1',\dots,A_{r-1}'$, where $A_i'$ has first term $a_i$ and common difference $d_i + 2nt$ are colour-focused at $f + (m-1)2n$, and also $\{f,f+2nt,f+2(2nt),\dots,f+(m-2)2nt$ is monochromatic of a different colour to the $A_i'$s. This gives $r$ colour-focused APs of length $m-1$.

    This completes the induction, which proves the claim, which completes the outer induction on $m$ and hence concludes the proof.
\end{proof}

\underline{Note}: for these proofs, focus on the pictures (even though I haven't drawn them here)!

\marginpar{Lecture 5}

\subsection*{The Ackermann/Grzegorczyk Hierarchy}

\begin{defin*} The \undf{Ackermann} or \undf{Grzegorcyzk} \undf{Hierarchy} is the sequence of funcgtions $f_1,f_2,\dots$ (each $\N\ra\N$) given by:
    \begin{itemize}
        \item $f_1(x) = 2x$
        \item $f_{n+1}(x) = f_n^{(x)}(1) = \underbrace{f_n(f_n(\dots f_n}_{x\textrm{ times}}(1)\dots))$
    \end{itemize}
\end{defin*}

Let's explore these a bit.
\begin{itemize}
    \item $f_2(x) = 2^x$
    \item $f_3(x) = \left. 2^{2^{.^{.^{2}}}} \right\rbrace x$
    \item $f_4(1) = 2$, $f_4(2) = 2^2 = 4$, $f_4(3) = 2^{2^{2^2}} = 65536$, $f_4(4) = \left. 2^{2^{.^{.^{2}}}} \right\rbrace 65536$ \it{etc}...
\end{itemize}

We say $f:\N\ra \N$ is of \undf{type $n$} if there exist $c,d > 0$ such that $f_n(cx) \le f(x) \le f_n(dx)$ for all $x$. So our upper bound for $W(3,k)$ was a function of $k$ of type 3 (note that even though it is a tower of $k$s, not $2$s, the height of the tower is far more significant). Our bound on $W(m,k)$ ($m$ fixed) is of type $m$.

In fact, our bound on $W(m) = W(m,2)$, as a function of $m$, grows faster than every $f_n$! [This means that $W(m)$ is \it{not} primitive recursive.] This is often a feature of such `double inductions', and for a long time it was thought that perhaps $W(m)$ really does grow this fast.

However, Shelah (1987) found a proof of VdW using induction only on $m$. His proof gives that $W(m,k) \le f_4(m+k)$. This isn't bad, but $f_4$ is still a pretty big function. Graham offered \$1000 for a bound on $W(m)$ that was $f_3(m)$. Gowers (1998) showed $W(m) \le 2^{2^{2^{2^{2^{m+9}}}}}$ - `almost' of type $2$ - a huge improvement on type 3.

What about a lower bound? It is known that $W(m) \ge \frac{2^m}{8m}$; this is comparatively extremely small.

\begin{cor}
    Whenever $\N$ is finitely coloured, some colour class arbitrarily long APs. \qedsymbol
\end{cor}
\begin{remark*}
    We \it{cannot} guarantee an infinite AP - \it{e.g.} R-BB-RRR-BBBB....
    
    Alternatively, list all infinite APs as $A_1,A_2,\dots$ (which we can do since they are countable). Pick distinct $x_1,y_1\in A_1$ and make $x_1$ red and $y_1$ blue. Now pick distinct $x_2,y_2 \in A_2 \backslash \{x_1,y_1\}$ and make $x_2$ red, $y_2$ blue. Continue. Then we wipe out all the APs.

    The idea here is that there simply aren't enough APs.
\end{remark*}

\begin{theorem}[Strengthened Van der Waerden Theorem]
     For all $m$, whenever $\N$ is finitely coloured there exists an AP of length $m$ that, together with its common difference, is monochromatic.
\end{theorem}
\begin{proof}
    We will prove this with VdW itself (this is the sign of a good theorem).

    We induct on $k$, the number of colours. $k = 1$ is trivial.

    Given $n$ suitable for $k-1$ (\it{i.e.} $n$ such that whenever $[n]$ is $k-1$-coloured, there exists a monochromatic AP+CD of length $m$), we'll show $W(n(m-1)+1,k)$ suitable for $k$.

    Given a $k$-colouring of $[W(n(m-1)+1,k)]$, we have a mono AP of length $n(m-1)+1$; say $\{a,a+d,\dots,a+2d,\dots,a+n(m-1)d\}$ is red.

    Now if $d$ is red then we are done by $\{a,a+d,\dots,a+(m-1)d\}\cup\{d\}$ is monochromatic. Similarly, if any $rd$, $1\le r \le n$ is red, then we are done since $\{a,a+rd,\dots, a + (m-1)rd\}\cup\{rd\}$ is monochromatic.

    Thus none of $\{d, 2d, \dots, nd\}$ is red, \it{i.e.} $\{d,2d,\dots,nd\}$ is $(k-1)$-coloured. So, by definition of $n$, we are done.
\end{proof}

What kind of bounds do we get here? Since we are iterating on the left hand side, this actually grows vastly more quickly than anything we have considered before; it's so mind bogglingly large that it essentially isn't worth thinking about.

\begin{remark*}[Remarks]\ 
    \begin{enumerate}
        \item From now on, we don't care aboud bounds.
        \item The case $m = 2$ is \it{Schur's Theorem}: Whenever $\N$ is finitely coloured, there exists monochromatic $x,y,z$ with $x + y = z$.
        \item Can also prove Schur direcly from Ramsey. Indeed, given a $k$-colouring $c$ of $\N$, induce a colouring of $[n]^{(2)}$ ($n$ large enough) by $d(ij) = c(j-i)$.
        
        By Ramsey, we can find a monochromatic triangle, say $ijk$. So $c(j-i) = c(k-j) = c(j-i)$, and $(j-i) + (k-j) = (k-i)$, so done.
    \end{enumerate}
\end{remark*}

\subsection{The Hales-Jewett Theorem}

You may have noticed that we haven't taken full advantage of the structure of the natural numbers - all we needed was some idea of equal spacing, and we didn't care so much about the additive structure. It turns out that we have a theorem that describes the same phenomenon but with all the clutter removed.

\begin{defin*}[Combinatorial Line]
Let $X$ be a finite set. A subset $L$ of $X^n$ (`the $n$-dimensional cube on alphabet $X$') is called a \undf{line} or \undf{combinatorial line} if there exists non-empty $I\subset [n]$ and $a_i \in X$, each $i \in [n] - I$ such that $$L = \{(x_1,\dots,x_n)\in X^n: x_i = a_i\ \forall i \not\in I, x_i = x_j\ \forall i,j\in I\}$$
We call $I$ the set of \undf{active coordinates}
\end{defin*}

For example, in $[3]^{(2)}$, lines can be:
\begin{itemize}
    \item vertical lines, where $I = \{2\}$
    \item horizontal lines, where $I = \{1\}$
    \item one diagonal line, where $I = \{1,2\}$ (the main diagonal only)
\end{itemize}
In $[3]^{3}$, we could have
\begin{itemize}
    \item $\{(1,2,1),(1,2,2),(1,2,3)\}$ with $I = \{3\}$, or
    \item $\{(1,3,1),(2,3,2),(3,3,3)\}$ with $I = \{1,3\}$, \it{etc}...
\end{itemize}

Note that the definition of `line' is invariant under reorderings of $x$ - so \it{e.g.} $\{1,3),(2,2),(3,1)\}$ is \it{not} a line in $[3]^{2}$.

\begin{theorem}[The Hales-Jewett Theorem]
    For all $m,k$ there exists $n$ such that whenever $[m]^{n}$ is $k$-coloured, there exists a monochromatic line.
\end{theorem}
\begin{remark*}[Remarks]\ 
    \begin{enumerate}
        \item The least such $n$ (if it exists) is denoted $\hj(m,k)$.
        \item A game of $m$-in-a-row Noughts \& Crosses, played in enough dimensions, cannot end in a draw! [Exercise: show that it is a first-player win (optional).]
        \item $\hj\implies$VdW immediately: just map $X^n$ linearly to $\N$ (note: perhaps requries positive coefficients).
        
        Indeed, given a $k$-colouring $c$ of $\N$, induce a $k$-colouring $d$ of $[m]^n$ ($n$ large) by $d((x_1,\dots,x_n)) = c(x_1+\dots+x_n)$. By assumption we have a monochromatic line $L$ in $[m]^{n}$, which corresponds to a monochromatic AP of length $m$ in $\N$ (with common difference = \# active coords of $L$).

        Thus we can view HJ as an `abstract version' of VdW.
    \end{enumerate}
\end{remark*}

Before we begin, we also need a few definitions:

For a line $L$ in $[m]^{(n)}$, write $L^-$ for its first point and $L^+$ for its last point (in the ordering on $[m]^{(n)}$) given by $x \le y$ if $x_i \le y_i$ for all $i$)

Say lines $L_1,\dots,L_r$ are \undf{focused} at $f$ if $L_i^+ = f$ for all $i$, and they are \undf{colour-focused} (for a given colouring) if in addition each $L_i - \{L_i^+\}$ is monochromatic, no two of the same colour.

\begin{proof}
   Induction on $m$. $m = 1$ is trivial.

   Given $m > 1$, we may assume $\hj(m-1,k)$ exists for all $k$.

   \underline{Claim}: for all $r\le k$, there exists an $n$ such that whenever $[m]^{(n)}$ is $k$-coloured, there exists either:
   \begin{itemize}
       \item a monochromatic line, or
       \item $r$ colour-focused lines
   \end{itemize}
   [Then done: put $r = k$ and look a the focus.]

   \underline{Proof of Claim}: Induction on $r$.

   $r = 1$ is done - take $n = \hj(m-1,k)$.

   Given $n$ suitable for $m-1$, we'll show $n + \hj(m-1,k^{m^n}) = : n + n'$ is suitable for $r$.

   So, given a $k$-colouring $c$ of $[m]^{n+n'}$ with no monochromatic line: View $[m]^{n+n'}$ as $[m]^{n}\times [m]^{n'}$. There are $k^{m^n}$ ways to colour $[m]^{n}\times [m]^{n'}$. At each point of $[m]^{n'}$ we have one of $k^{m^n}$ ``patterns''.

   So by choice of $n'$, we have a line $L$ in $[m]^{n'}$ such that for all $a \in [m]^{n}$, for all $b,b' \in L - \{L^+\}$ we have $c(a,b) = c(a,b') =: c'(a)$, say. By definition of $n$, $c'$ has $r-1$ colour-focused lines, say $L_1,\dots,L_{r-1}$ with active coordinate sets $J_1,\dots,J_{r-1}$ respectively, focused at $f$. Let $L$ have active coord set $I$. Then the lines $L_1',\dots,L_{r-1}'$, where $L_i'$ starts at $(L_i^-,L^-)$ and has active coord set $J_i\cup I$, are colour-focused at $(f,L^+)$. Also the line starting $(f,L^-)$ with active coord set $I$ is monochromatic (apart from final point), of a different colour to the $L_i'$s, giving $r$ colour-focused lines.

   This completes the induction, hence the proof of the claim, has the whole proof.
\end{proof}

What is a line though? If you think about it, it's just a one-dimensional subspace. This begs the question, can we perhaps get a monochromatic \it{two}-dimesional subspace?

A \undf{$d$-parameter set} or \undf{$d$-dimensional subspace} of $X^n$ is a set $S \subset X^n$ such that there exist disjoint, non-empty $I_1,\dots,I_d\subset [n]$ and $a_i \in X$, each $i \in [n]-(I_1\cup\dots\cup I_d)$ such that
\begin{align*}
    S = \{x \in X^n : x_i = a_i\ \forall i\not\in I_1\cup\dots\cup I_d,\textrm{ and }x_i = x_j\ \forall i,j\textrm{ with }i,j\in I_k\textrm{ for some }k\}
\end{align*}
For instance, in $X^3$:
\begin{itemize}
    \item $\{(x,y,2):x,y\in X\}$ is a $2$-parameter set: $I_1 = \{1\}$, $I_2 = \{2\}$
    \item $\{(x,y,x):x,y\in X\}$ is a $2$-parameter set: $I_1 = \{3\}$, $I_2 = \{2\}$
\end{itemize}

This is in fact true.

\begin{theorem}[Extended Hales-Jewett Theorem]
    For all $m,k,d$ ther exists $m$ such that whenever $[m]^{n}$ is $k$-coloured, there exists a monochromatic $d$-parameter set.
\end{theorem}
\begin{proof}
    View $X^{dn}$ (the $dn$-dimensional cube over alphabet $X$) as $(X^d)^n$, the $n$-dimensional cube on alphabet $X^d$. Clearly any line in $(X^d)^n$ (alphabet $X^d$) corresponds to a $d$-parameter set in $X^{dn}$ (alphabet $X$). So we are done - we can take $n = d\cdot \hj(m^d,k)$.
\end{proof}

Let $S$ be a finite subset of $\N^d$. A \undf{homothetic copy} of $A$ is any set of the form $a + \lambda S$, where $a \in \N^d$ and $\lambda \in S$.

For instance, in $\N^1$, a homothetic copy of $\{1,2,\dots,m\}$ is precisely an AP of length $m$.

Can we guarantee a homothetic copy of $S$ that is monochromatic? For instance, in $\N^2$, can we find a monochromatic square?

\begin{theorem}[Gallai's Theorem]
    Let $S$ be a finite subset of $\N^d$. Then, whenever $\N^d$ is fiinitely coloured, there exists a monochromatic homothetic copy of $S$.
\end{theorem}
\begin{proof}
    Let $S = \{S(1),\dots,S(m)\}$. Given a $k$-colouring $c$ of $\N^d$, induce a $k$-colouring $c'$ of $[m]^n$ ($n$ large) by: $c'(x_1,x_2,\dots,x_d) = c(S(x_1)+\dots + S(x_d))$.

    Then we have a monochromatic line $L$ for $c'$ ($n$ large), say with active coord set $I$. But now, taking $S(x_1)+\dots + S(x_d)$ for each $x \in L$, we have a monochromatic, homothetic copy of $S$ (with $\lambda = |I|$).
\end{proof}
\begin{remark*}[Remarks]\ 
    \begin{enumerate}
        \item Suppose $S = \{(0,0),(1,0),(0,1),(1,1)\} \subset \N^2$. Then we get a monochromatic square. Could we instead have applied Extended HJ, $d = 2$, on alphabet of size $2$?
        
        The answer is \it{no}: this would only give a monochromatic rectangle.

        \item Or, can prove Gallai by product arguments and focusing (don't try it, symbol overload, but similar argument).
    \end{enumerate}
\end{remark*}

\marginpar{Lecture 7}

\end{document}